\newpage
\chapter{Kinematic Constraints and Rigid Couplings}\label{chap:kinematic-constraints}

Finite element models frequently need relations between degrees of freedom that
are not represented by the ordinary structural elements themselves. A flange may
be treated as rigid, a force may be introduced through a point outside the mesh,
a bearing may be represented by a bushing acting on an entire hole, or two nodes
may be required to remain at a prescribed distance. These modelling operations
look different in a graphical user interface, but they are all examples of
\textit{kinematic constraints}: relations that prescribe which relative motions
are admissible.

The mathematical machinery for imposing such relations was developed in
Chapter~\ref{chap:multifreedom-constraints}. The present chapter therefore does
not repeat the penalty, elimination, Lagrange multiplier and augmented Lagrangian
methods. Instead, it asks a different question: \textit{what constraint should be
constructed for a particular engineering purpose?}

This distinction is important. The same engineering constraint may be enforced
by different numerical techniques, and two solver commands that look similar may
produce different algebraic systems. Conversely, commands with very different
names may generate essentially the same constraint equations. The familiar names
RBE1, RBE2 and RBE3 are consequently used in this book as general engineering
names for three widely recognized classes of coupling, not as an attempt to
follow the terminology of a particular finite element program.

Contact provides the complementary case. Chapter~\ref{chap:contact-formulations}
showed that ordinary contact is not a permanent equality relation but a unilateral
constraint whose active set is part of the solution. The distinction can be
summarized as

\begin{equation}
    \begin{aligned}
        \text{bilateral kinematic constraint:} && \m{g}(\m{u}) &= \m{0}, \\
        \text{unilateral contact constraint:} && g_n(\m{u}) &\geq 0,
    \end{aligned}
\end{equation}

with the contact reaction additionally governed by the complementarity conditions
introduced in Chapter~\ref{chap:contact-formulations}.

\section{From General Constraints to Engineering Couplings}

A general linear multifreedom constraint may be written in the form used in
Chapter~\ref{chap:multifreedom-constraints},

\begin{equation}\label{kinematic-general-linear-constraint}
    \m{A}\m{u}=\m{b}.
\end{equation}

The coefficients of $\m{A}$ may be entered explicitly, but doing so becomes
cumbersome when the relation is dictated by geometry. A rigid spider with one
reference point and fifty attachment nodes already requires hundreds of
coefficients. Repeating the construction for each connection is tedious and,
more importantly, easy to get wrong.

Rigid and interpolation couplings can therefore be understood as \textit{constraint
generators}. The analyst specifies geometry, active degrees of freedom and,
where necessary, weights; the program constructs the corresponding equations.
The generated relation can then be imposed by one of the numerical procedures
discussed in Chapter~\ref{chap:multifreedom-constraints}.

This viewpoint also clarifies why a coupling should not automatically be regarded
as a physical finite element. An ideal RBE2 or RBE3 usually has no constitutive
law, no Young's modulus and no strain energy of its own. It changes the set of
admissible displacement fields. A beam, spring, bushing or adhesive layer is
fundamentally different because it stores energy according to a constitutive law.

\begin{bbox}
\textbf{Engineering note -- Constraint definition and constraint enforcement are different questions.}

A fixed-distance relation, an RBE2 rigid coupling and an RBE3 interpolation
coupling describe different admissible motions. Each may in principle be imposed
by elimination, multipliers, a penalty-like method or another solver-specific
algorithm. The physical definition should therefore be selected first; the
numerical enforcement method is a second decision.
\end{bbox}

\section{Evolution of the Rigid-Element Family}

Early finite element programs exposed general multipoint or constraint equations
directly. This remains the most general representation: any linear relation that
can be written as Eq.~\eqref{kinematic-general-linear-constraint} can be entered
coefficient by coefficient. Recurrent engineering patterns, however, soon led to
specialized definitions.

The RBE family is best understood as an evolution of these recurring patterns:

\begin{enumerate}
    \item \textbf{RBE1} provides a general rigid coupling between selected sets
          of independent and dependent degrees of freedom. It is close in spirit
          to a structured multipoint-constraint definition.
    \item \textbf{RBE2} expresses the dependent-node motion directly from the
          translation and rotation of a reference node. It therefore creates an
          exactly rigid kinematic region in the selected degrees of freedom.
    \item \textbf{RBE3} was introduced for the opposite modelling need: transfer
          forces, moments and generalized motion through a set of surrounding
          nodes \textit{without making that surrounding region rigid}. Its
          reference motion is obtained by interpolation, commonly by a weighted
          least-squares construction.
\end{enumerate}

Modern programs frequently place another abstraction above these definitions:
remote points, reference points, surface-based constraints, joints and coupling
objects. Such objects may internally generate classical constraint equations,
interpolation relations, contact-based MPCs or dedicated constraint elements.
The GUI object is consequently not the theory. The engineering behaviour of the
underlying relation is what matters.

\section{Classification of Kinematic Constraints}

It is useful to classify a coupling by the relative motion it permits before
considering any solver command.

\begin{center}
\begin{tabular}{|p{0.22\textwidth}|p{0.33\textwidth}|p{0.32\textwidth}|}
\hline
\textbf{Constraint type} & \textbf{Prescribed relation} & \textbf{Typical use} \\
\hline
Equal DOF & Selected degrees of freedom have the same value & Coincident-node tying, symmetry, simple joints \\
\hline
General linear MPC & Arbitrary linear combination of DOFs & Periodicity, special kinematic relations \\
\hline
RBE1 & Selected rigid relations between sets of DOFs & General rigid coupling \\
\hline
RBE2 & Rigid-body motion about a reference node & Rigid flange, rigid spider, ideal rigid attachment \\
\hline
RBE3 & Weighted interpolation of reference motion & Load distribution, remote support, deformable attachment \\
\hline
Fixed distance & Constant distance between geometric points & Links, rods, mechanism constraints \\
\hline
Fixed angle / joint & Selected relative rotations or translations suppressed & Revolute, cylindrical, spherical and other joints \\
\hline
Contact & Inequality and complementarity conditions & Separation, sliding, impact and friction \\
\hline
\end{tabular}
\end{center}

The last row is deliberately included even though contact was treated in the
preceding chapter. It emphasizes that contact and permanent couplings answer the
same broad engineering question -- which relative motion is admissible -- but
have fundamentally different mathematical status.

\section{Mathematical Background}

For a bilateral nonlinear constraint the most general notation used in this
chapter is

\begin{equation}\label{kinematic-general-nonlinear-constraint}
    \m{g}(\m{u})=\m{0}.
\end{equation}

Linearization around the current iterate $\m{u}^{(k)}$ gives

\begin{equation}
    \m{g}(\m{u}^{(k)})+
    \m{G}^{(k)}\Delta\m{u}=\m{0},
    \qquad
    \m{G}^{(k)}=\left.\frac{\partial\m{g}}{\partial\m{u}}\right|_{\m{u}^{(k)}}.
\end{equation}

If the relation is linear and fixed in the initial configuration,
$\m{G}=\m{A}$ is constant and Eq.~\eqref{kinematic-general-linear-constraint}
is recovered. This apparently small distinction becomes crucial for large
rotations. A constraint constructed once from the initial geometry follows the
initial tangent; a geometrically exact constraint must update its direction and
moment arms as the configuration changes.

When the constraint is enforced by Lagrange multipliers, for example, the
linearized equilibrium system has the familiar saddle-point form

\begin{equation}\label{kinematic-lagrange-system}
    \begin{bmatrix}
        \m{K}_t & \m{G}^T \\
        \m{G}   & \m{0}
    \end{bmatrix}
    \begin{bmatrix}
        \Delta\m{u} \\
        \Delta\m{\lambda}
    \end{bmatrix}
    =
    \begin{bmatrix}
        \m{r} \\
        -\m{g}
    \end{bmatrix}.
\end{equation}

The multipliers $\m{\lambda}$ have the physical interpretation of constraint
reactions. Equation~\eqref{kinematic-lagrange-system} is included here only to
show the distinction between the \textit{definition} $\m{g}$ and its
\textit{enforcement}; the numerical properties of multiplier systems were
discussed in Chapter~\ref{chap:multifreedom-constraints}.

\section{RBE1 -- General Rigid Coupling}

RBE1 is the least specialized member of the classical family. It is useful to
think of it as a structured way of defining a rigid relation between selected
degrees of freedom rather than as a spider with one mandatory reference point.
A set of independent degrees of freedom supplies the retained motion and a set of
dependent degrees of freedom is related to it through rigid-body compatibility.

Its historical importance is greater than its present visibility. Modern
preprocessors often expose more convenient RBE2, RBE3, joint or remote-coupling
objects, while a general constraint-equation facility covers cases that once
motivated RBE1 directly. Nevertheless, RBE1 remains conceptually useful because
it sits between a completely arbitrary MPC and the highly structured RBE2.

The important modelling question is not whether the solver literally contains an
object named RBE1, but whether the generated equations impose a rigid relation
between the selected DOFs. If they do, the same cautions concerning redundant
constraints, thermal expansion and finite rotations apply.

\begin{bbox}
\textbf{Common mistake -- Treating the RBE name as the formulation.}

Commercial programs expose the same engineering relation under different names.
The command name should never replace checking which DOFs are dependent, whether
the surrounding region is made rigid, how the equations are enforced and whether
the formulation is valid for the selected analysis type.
\end{bbox}

\section{RBE2 -- Rigid-Body Coupling}

\begin{bbox}[0.95]
    The rigid freebody relations between \textbf{slave} and \textbf{master} node are:
    \begin{eqarray}
        \m{T}_s &= \m{T}_m + \m{R}_m \times \overline{\m{x}} \\
        \m{R}_s &= \m{R}_m
    \end{eqarray}

    Where $ \overline{\m{x}} $ is a vector from \textbf{slave} to \textbf{master}:
    \begin{equation}
        \overline{\m{x}} =
        \begin{bmatrix}
            \overline{x} \\
            \overline{y} \\
            \overline{z}
        \end{bmatrix} =
        \begin{bmatrix}
            x_s - x_m \\
            y_s - y_m \\
            z_s - z_m
        \end{bmatrix}
    \end{equation}

    Then \textit{vector multiplication} (cross product) is defined as:
    \begin{eqarray}
        \m{R}_s \times \m{x}_s
        &= \m{R}_m \times \m{x} = & \\
        &= \begin{bmatrix}
             \varphi_m \\
             \psi_m \\
             \theta_m
           \end{bmatrix} \times \begin{bmatrix}
             x_s - x_m \\
             y_s - y_m \\
             z_s - z_m
           \end{bmatrix}
        &= \begin{bmatrix}
             \varphi_m \\
             \psi_m \\
             \theta_m
           \end{bmatrix} \times \begin{bmatrix}
             \overline{x} \\
             \overline{y} \\
             \overline{z}
           \end{bmatrix} = \\
        &= \begin{bmatrix}
             \psi_m (z_s - z_m ) - \theta_m (y_s - y_m) \\
             \theta_m (x_s - x_m) - \varphi_m (z_s - z_m) \\
             \varphi_m (y_s - y_m) - \psi_m (x_s - x_m)
           \end{bmatrix}
        &= \begin{bmatrix}
            \psi_m \overline{z} - \theta_m \overline{y} \\
            \theta_m \overline{x} - \varphi_m \overline{z} \\
            \varphi_m \overline{y} - \psi_m \overline{x}
        \end{bmatrix} = \\
        &= \begin{bmatrix}
            \phantom{-}0 & \phantom{-}\overline{z} & -\overline{y} \\
            -\overline{z} & \phantom{-}0 & \phantom{-}\overline{x} \\
            \phantom{-}\overline{y} & -\overline{x} & \phantom{-}0
        \end{bmatrix}
        \begin{bmatrix}
            \varphi_m \\
            \psi_m \\
            \theta_m
        \end{bmatrix}
    \end{eqarray}
\end{bbox}

\begin{bbox}[0.95]
    The \textbf{slave} equations can be written as:
    \begin{equation}
        \begin{bmatrix}
            u_s \\
            v_s \\
            w_s \\
            \varphi_s \\
            \psi_s \\
            \theta_s
        \end{bmatrix}
        = \begin{bmatrix}
            1 & 0 & 0 & 0 & \phantom{-(}z_s - z_m & -(y_s - y_m) \\
            0 & 1 & 0 & -(z_s - z_m) & 0 & \phantom{-(}x_s - x_m \\
            0 & 0 & 1 & \phantom{-(}y_s - y_m & -(x_s - x_m) & 0 \\
            0 & 0 & 0 & 1 & 0 & 0 \\
            0 & 0 & 0 & 0 & 1 & 0 \\
            0 & 0 & 0 & 0 & 0 & 1 \\
        \end{bmatrix}
        \begin{bmatrix}
            u_m \\
            v_m \\
            w_m \\
            \varphi_m \\
            \psi_m \\
            \theta_m
        \end{bmatrix}
    \end{equation}

    or:
    \begin{equation}
        \begin{bmatrix}
            u_s \\
            v_s \\
            w_s \\
            \varphi_s \\
            \psi_s \\
            \theta_s
        \end{bmatrix}
        = \begin{bmatrix}
            1 & 0 & 0 & \phantom{-}0 & \phantom{-}\overline{z} & -\overline{y} \\
            0 & 1 & 0 & -\overline{z} & \phantom{-}0 & \phantom{-}\overline{x} \\
            0 & 0 & 1 & \phantom{-}\overline{y} & -\overline{x} & \phantom{-}0 \\
            0 & 0 & 0 & \phantom{-}1 & \phantom{-}0 & \phantom{-}0 \\
            0 & 0 & 0 & \phantom{-}0 & \phantom{-}1 & \phantom{-}0 \\
            0 & 0 & 0 & \phantom{-}0 & \phantom{-}0 & \phantom{-}1 \\
        \end{bmatrix}
        \begin{bmatrix}
            u_m \\
            v_m \\
            w_m \\
            \varphi_m \\
            \psi_m \\
            \theta_m
        \end{bmatrix}
    \end{equation}
\end{bbox}

\newpage
\subsection{RBE2 Coefficients Example:}

Master Node:
\begin{equation}
    \m{X}_m =
    \begin{bmatrix}
        1 & 1 & 1
    \end{bmatrix}^T
\end{equation}

Slave Nodes:
\begin{eqarray}
    \m{X}_{s,1} =
    \begin{bmatrix}
        0 & 0 & 0
    \end{bmatrix}^T \\
    \m{X}_{s,1} =
    \begin{bmatrix}
        2 & 0 & 0
    \end{bmatrix}^T \\
    \m{X}_{s,1} =
    \begin{bmatrix}
        2 & 2 & 0
    \end{bmatrix}^T \\
    \m{X}_{s,1} =
    \begin{bmatrix}
        0 & 2 & 0
    \end{bmatrix}^T \\
\end{eqarray}

Then Slave Node 1 equations are:
\begin{eqarray}
    \m{x}_1 &=
    \begin{bmatrix}
        x_{s,1} - x_m & y_{s,1} - y_m & z_{s,1} - z_m
    \end{bmatrix}^T =
    \begin{bmatrix}
        -1 & -1 & -1
    \end{bmatrix}^T
\end{eqarray}

Other Slave Nodes:
\begin{eqarray}
    \m{x}_2 &=
    \begin{bmatrix}
        2 - 1 & 0 - 1 & 0 - 1
    \end{bmatrix}^T =
    \begin{bmatrix}
        \phantom{-}1 & -1 & -1
    \end{bmatrix}^T \\
    \m{x}_3 &=
    \begin{bmatrix}
        2 - 1 & 2 - 1 & 0 - 1
    \end{bmatrix}^T =
    \begin{bmatrix}
        \phantom{-}1 & \phantom{-}1 & -1
    \end{bmatrix}^T \\
    \m{x}_4 &=
    \begin{bmatrix}
        0 - 1 & 2 - 1 & 0 - 1
    \end{bmatrix}^T =
    \begin{bmatrix}
        -1 & \phantom{-}1 & -1
    \end{bmatrix}^T
\end{eqarray}

Therefore, after filling in the coefficients:
\renewcommand\arraystretch{1.6}
\begin{eqarray}
    \begin{bmatrix}
        u_{s,1} \\
        v_{s,1} \\
        w_{s,1} \\
        \varphi_{s,1} \\
        \psi_{s,1} \\
        \theta_{s,1}
    \end{bmatrix}
    &= \begin{bmatrix}
        1 & 0 & 0 & \phantom{-}0 & -1 & \phantom{-}1 \\
        0 & 1 & 0 & \phantom{-}1 & \phantom{-}0 & -1 \\
        0 & 0 & 1 & -1 & \phantom{-}1 & \phantom{-}0 \\
        0 & 0 & 0 & \phantom{-}1 & \phantom{-}0 & \phantom{-}0 \\
        0 & 0 & 0 & \phantom{-}0 & \phantom{-}1 & \phantom{-}0 \\
        0 & 0 & 0 & \phantom{-}0 & \phantom{-}0 & \phantom{-}1 \\
    \end{bmatrix}
    \begin{bmatrix}
        u_m \\
        v_m \\
        w_m \\
        \varphi_m \\
        \psi_m \\
        \theta_m
    \end{bmatrix} \\
    \begin{bmatrix}
        u_{s,2} \\
        v_{s,2} \\
        w_{s,2} \\
        \varphi_{s,2} \\
        \psi_{s,2} \\
        \theta_{s,2}
    \end{bmatrix}
    &= \begin{bmatrix}
        1 & 0 & 0 & \phantom{-}0 & -1 & \phantom{-}1 \\
        0 & 1 & 0 & \phantom{-}1 & \phantom{-}0 & \phantom{-}1 \\
        0 & 0 & 1 & -1 & -1 & \phantom{-}0 \\
        0 & 0 & 0 & \phantom{-}1 & \phantom{-}0 & \phantom{-}0 \\
        0 & 0 & 0 & \phantom{-}0 & \phantom{-}1 & \phantom{-}0 \\
        0 & 0 & 0 & \phantom{-}0 & \phantom{-}0 & \phantom{-}1 \\
    \end{bmatrix}
    \begin{bmatrix}
        u_m \\
        v_m \\
        w_m \\
        \varphi_m \\
        \psi_m \\
        \theta_m
    \end{bmatrix} \\
    \begin{bmatrix}
        u_{s,3} \\
        v_{s,3} \\
        w_{s,3} \\
        \varphi_{s,3} \\
        \psi_{s,3} \\
        \theta_{s,3}
    \end{bmatrix}
    &= \begin{bmatrix}
        1 & 0 & 0 & \phantom{-}0 & -1 & -1 \\
        0 & 1 & 0 & \phantom{-}1 & \phantom{-}0 & \phantom{-}1 \\
        0 & 0 & 1 & \phantom{-}1 & -1 & \phantom{-}0 \\
        0 & 0 & 0 & \phantom{-}1 & \phantom{-}0 & \phantom{-}0 \\
        0 & 0 & 0 & \phantom{-}0 & \phantom{-}1 & \phantom{-}0 \\
        0 & 0 & 0 & \phantom{-}0 & \phantom{-}0 & \phantom{-}1 \\
    \end{bmatrix}
    \begin{bmatrix}
        u_m \\
        v_m \\
        w_m \\
        \varphi_m \\
        \psi_m \\
        \theta_m
    \end{bmatrix} \\
    \begin{bmatrix}
        u_{s,4} \\
        v_{s,4} \\
        w_{s,4} \\
        \varphi_{s,4} \\
        \psi_{s,4} \\
        \theta_{s,4}
    \end{bmatrix}
    &= \begin{bmatrix}
        1 & 0 & 0 & \phantom{-}0 & -1 & -1 \\
        0 & 1 & 0 & \phantom{-}1 & \phantom{-}0 & -1 \\
        0 & 0 & 1 & \phantom{-}1 & \phantom{-}1 & \phantom{-}0 \\
        0 & 0 & 0 & \phantom{-}1 & \phantom{-}0 & \phantom{-}0 \\
        0 & 0 & 0 & \phantom{-}0 & \phantom{-}1 & \phantom{-}0 \\
        0 & 0 & 0 & \phantom{-}0 & \phantom{-}0 & \phantom{-}1 \\
    \end{bmatrix}
    \begin{bmatrix}
        u_m \\
        v_m \\
        w_m \\
        \varphi_m \\
        \psi_m \\
        \theta_m
    \end{bmatrix}
\end{eqarray}
\renewcommand\arraystretch{1}

When all put together:
\begin{equation}
    \begin{bmatrix}
        u_{s,1} \\
        v_{s,1} \\
        w_{s,1} \\
        \varphi_{s,1} \\
        \psi_{s,1} \\
        \theta_{s,1} \\
        u_{s,2} \\
        v_{s,2} \\
        w_{s,2} \\
        \varphi_{s,2} \\
        \psi_{s,2} \\
        \theta_{s,2} \\
        u_{s,3} \\
        v_{s,3} \\
        w_{s,3} \\
        \varphi_{s,3} \\
        \psi_{s,3} \\
        \theta_{s,3} \\
        u_{s,4} \\
        v_{s,4} \\
        w_{s,4} \\
        \varphi_{s,4} \\
        \psi_{s,4} \\
        \theta_{s,4}
    \end{bmatrix}
    = \begin{bmatrix}
        1 & 0 & 0 & \phantom{-}0 & -1 & \phantom{-}1 \\
        0 & 1 & 0 & \phantom{-}1 & \phantom{-}0 & -1 \\
        0 & 0 & 1 & -1 & \phantom{-}1 & \phantom{-}0 \\
        0 & 0 & 0 & \phantom{-}1 & \phantom{-}0 & \phantom{-}0 \\
        0 & 0 & 0 & \phantom{-}0 & \phantom{-}1 & \phantom{-}0 \\
        0 & 0 & 0 & \phantom{-}0 & \phantom{-}0 & \phantom{-}1 \\
        1 & 0 & 0 & \phantom{-}0 & -1 & \phantom{-}1 \\
        0 & 1 & 0 & \phantom{-}1 & \phantom{-}0 & \phantom{-}1 \\
        0 & 0 & 1 & -1 & -1 & \phantom{-}0 \\
        0 & 0 & 0 & \phantom{-}1 & \phantom{-}0 & \phantom{-}0 \\
        0 & 0 & 0 & \phantom{-}0 & \phantom{-}1 & \phantom{-}0 \\
        0 & 0 & 0 & \phantom{-}0 & \phantom{-}0 & \phantom{-}1 \\
        1 & 0 & 0 & \phantom{-}0 & -1 & -1 \\
        0 & 1 & 0 & \phantom{-}1 & \phantom{-}0 & \phantom{-}1 \\
        0 & 0 & 1 & \phantom{-}1 & -1 & \phantom{-}0 \\
        0 & 0 & 0 & \phantom{-}1 & \phantom{-}0 & \phantom{-}0 \\
        0 & 0 & 0 & \phantom{-}0 & \phantom{-}1 & \phantom{-}0 \\
        0 & 0 & 0 & \phantom{-}0 & \phantom{-}0 & \phantom{-}1 \\
        1 & 0 & 0 & \phantom{-}0 & -1 & -1 \\
        0 & 1 & 0 & \phantom{-}1 & \phantom{-}0 & -1 \\
        0 & 0 & 1 & \phantom{-}1 & \phantom{-}1 & \phantom{-}0 \\
        0 & 0 & 0 & \phantom{-}1 & \phantom{-}0 & \phantom{-}0 \\
        0 & 0 & 0 & \phantom{-}0 & \phantom{-}1 & \phantom{-}0 \\
        0 & 0 & 0 & \phantom{-}0 & \phantom{-}0 & \phantom{-}1 \\
    \end{bmatrix}
    \begin{bmatrix}
        u_m \\
        v_m \\
        w_m \\
        \varphi_m \\
        \psi_m \\
        \theta_m
    \end{bmatrix}
\end{equation}

It can be read as:
\begin{eqarray}
    u_{s,1} &= u_m -\psi_m + \theta_m \\
    v_{s,1} &= v_m + \varphi_m - \theta_m \\
    \vdots \\
    \theta_{s,4} &= \theta_m
\end{eqarray}

Or:
\begin{equation}\label{rbe2-coefficients}
    \m{u}_s = \m{\closure{T}} \m{u}_m
\end{equation}


\subsection{Editing the Master Stiffness Matrix}

This equation should be then expanded to all DOFs, where in $ \m{T} $ the
slave DOFs columns are missing and the DOFs that do not take part in RBEs have $ 1 $ on
the diagonal. The resulting  matrix should have the same number of rows as
the master stiffness matrix.

Then the master stiffness equation is modified:
\begin{eqarray}\label{rbe2-master-slave-modified}
    \m{\hat{K}} &= \m{T}^T \m{K} \m{T} \\
    \m{\hat{f}} &= \m{T}^T \m{f} \\
    \m{\hat{K}} \m{\hat{u}} &= \m{\hat{f}} \\
\end{eqarray}

If we'd like to make the procedure more general, partition the master stiffness
equations such that:
\begin{equation}
    \begin{bmatrix}
        \m{K}_{uu} & \m{K}_{um} & \m{K}_{us} \\
        \m{K}_{um}^T & \m{K}_{mm} & \m{K}_{ms} \\
        \m{K}_{us}^T & \m{K}_{ms}^T & \m{K}_{ss}
    \end{bmatrix}
    \begin{bmatrix}
        \m{u}_u \\
        \m{u}_m \\
        \m{u}_s
    \end{bmatrix}
    = \begin{bmatrix}
        \m{f}_u \\
        \m{f}_m \\
        \m{f}_s
    \end{bmatrix}
\end{equation}

Based on \eqref{rbe2-master-slave-modified} we could probably partition
the matrices as:
\begin{equation}\label{rbe2-partition-T-u-m-s}
    \begin{matrix}
        \begin{matrix}
            \begin{matrix}
                \ 
            \end{matrix} \\
            \left. \begin{matrix}
                u \\
                m \\
                s \\
            \end{matrix} \right|
        \end{matrix}
        & \begin{matrix}
            \begin{matrix}
                u & m \\
                \hline
            \end{matrix} \\
            \begin{bmatrix}
                \m{I} & \m{0} \\
                \m{0} & \m{I} \\
                \m{0} & \m{\closure{T}} \\
            \end{bmatrix} \\
        \end{matrix}
    \end{matrix}
\end{equation}

If we expand \eqref{rbe2-coefficients} and partition the matrices
based on \eqref{rbe2-partition-T-u-m-s} into independent or
\textbf{unconstrained}, \textbf{master}
and \textbf{slave} nodes, then the equation:
\begin{equation}
    \m{u} = \m{T} \m{\hat{u}}
\end{equation}

becomes:
\begin{equation}
    \begin{bmatrix}
        \m{u}_u  \\
        \m{u}_m \\
        \m{u}_s \\
    \end{bmatrix}
    = \begin{bmatrix}
        \m{I} & \m{0} \\
        \m{0} & \m{I} \\
        \m{0} & \m{\closure{T}} \\
    \end{bmatrix}
    \begin{bmatrix}
        \m{u}_u \\
        \m{u}_m \\
    \end{bmatrix}
\end{equation}


Finally the equation \eqref{rbe2-master-slave-modified} can be rewritten as:
\begin{equation}
    \begin{bmatrix}
        \m{I} & \m{0} & \m{0} \\
        \m{0} & \m{I} & \m{\closure{T}}^T \\
    \end{bmatrix}
    \begin{bmatrix}
        \m{K}_{uu} & \m{K}_{um} & \m{K}_{us} \\
        \m{K}_{mu} & \m{K}_{mm} & \m{K}_{ms} \\
        \m{K}_{su} & \m{K}_{sm} & \m{K}_{ss} \\
    \end{bmatrix}
    \begin{bmatrix}
        \m{I} & \m{0} \\
        \m{0} & \m{I} \\
        \m{0} & \m{\closure{T}} \\
    \end{bmatrix}
    \begin{bmatrix}
        \m{u}_u \\
        \m{u}_m \\
    \end{bmatrix}
    = \begin{bmatrix}
        \m{I} & \m{0} & \m{0} \\
        \m{0} & \m{I} & \m{\closure{T}}^T \\
    \end{bmatrix}
    \begin{bmatrix}
        \m{f}_u \\
        \m{f}_m \\
        \m{f}_s \\
    \end{bmatrix}
\end{equation}

Expanding and multiplying:
\begin{equation}
    \scalebox{0.92}{
    $ \begin{bmatrix}
        \m{K}_{uu} & \m{K}_{um} & \m{K}_{us} \\
        \m{K}_{mu} + \m{\closure{T}}^T \m{K}_{su} &
        \m{K}_{mm} + \m{\closure{T}}^T \m{K}_{sm} &
        \m{K}_{ms} + \m{\closure{T}}^T \m{K}_{ss} \\
    \end{bmatrix}
    \begin{bmatrix}
        \m{I} & \m{0} \\
        \m{0} & \m{I} \\
        \m{0} & \m{\closure{T}} \\
    \end{bmatrix}
    \begin{bmatrix}
        \m{u}_u \\
        \m{u}_m \\
    \end{bmatrix}
    = \begin{bmatrix}
        \m{f}_u \\
        \m{f}_m + \m{\closure{T}}^T \m{f}_s \\
    \end{bmatrix} $
    }
\end{equation}

\begin{equation}
    \scalebox{0.95}{
    $ \begin{bmatrix}
        \m{K}_{uu} & \m{K}_{um} + \m{K}_{us} \m{\closure{T}} \\
        \m{K}_{mu} + \m{\closure{T}}^T \m{K}_{su} &
        \m{K}_{mm} + \m{\closure{T}}^T \m{K}_{sm} +
        \m{K}_{ms} \m{\closure{T}} + \m{\closure{T}}^T \m{K}_{ss} \m{\closure{T}} \\
    \end{bmatrix}
    \begin{bmatrix}
        \m{u}_u \\
        \m{u}_m \\
    \end{bmatrix}
    = \begin{bmatrix}
        \m{f}_u \\
        \m{f}_m + \m{\closure{T}}^T \m{f}_s \\
    \end{bmatrix} $
    }
\end{equation}

which is equal to:
\begin{equation}
    \scalebox{0.95}{
    $ \begin{bmatrix}
        \m{K}_{uu} & \m{K}_{um} + \m{K}_{us} \m{\closure{T}} \\
        (\m{K}_{um} + \m{K}_{us} \m{\closure{T}})^T &
        \m{K}_{mm} + \m{\closure{T}}^T \m{K}_{ms}^T +
        \m{K}_{ms} \m{\closure{T}} + \m{\closure{T}}^T \m{K}_{ss} \m{\closure{T}} \\
    \end{bmatrix}
    \begin{bmatrix}
        \m{u}_u \\
        \m{u}_m \\
    \end{bmatrix}
    = \begin{bmatrix}
        \m{f}_u \\
        \m{f}_m + \m{\closure{T}}^T \m{f}_s \\
    \end{bmatrix} $
    }
\end{equation}

considering that $ \m{f}_s = \m{0} $ the equations reduce to:
\begin{equation}
    \begin{bmatrix}
        \m{K}_{uu} & \m{\hat{K}}_{um} \\
        \m{\hat{K}}_{um}^T & \m{\hat{K}}_{mm} \\
    \end{bmatrix}
    \begin{bmatrix}
        \m{u}_u \\
        \m{u}_m \\
    \end{bmatrix}
    = \begin{bmatrix}
        \m{f}_u \\
        \m{f}_m \\
    \end{bmatrix}
\end{equation}

where:
\begin{eqarray}
    \m{\hat{K}}_{um} &= \m{K}_{um} + \m{K}_{us} \m{\closure{T}} \\
    \m{\hat{K}}_{mm} &= \m{K}_{mm} + \m{\closure{T}}^T \m{K}_{ms}^T
                      + \m{K}_{ms} \m{\closure{T}}
                      + \m{\closure{T}}^T \m{K}_{ss} \m{\closure{T}} \\
                     &= \m{K}_{mm} + \m{K}_{ms} \m{\closure{T}}
                      + (\m{K}_{ms} \m{\closure{T}})^T
                      + \m{\closure{T}}^T \m{K}_{ss} \m{\closure{T}}
\end{eqarray}


\begin{bbox}

    Then if \textbf{master-slave} relations are written as (\textit{general case}):
    \begin{equation}
        \m{u}_s = \m{\closure{T}} \m{u}_m + \m{g}
    \end{equation}

    Inserting into the partitioned master stiffness equations:
    \begin{equation}
        \begin{bmatrix}
            \m{K}_{uu} & \m{\hat{K}}_{um} \\
            \m{\hat{K}}_{um}^T & \m{\hat{K}}_{mm}
        \end{bmatrix}
        \begin{bmatrix}
            \m{u}_u \\
            \m{u}_m
        \end{bmatrix}
        = \begin{bmatrix}
            \m{f}_u - \m{K}_{us} \m{g} \\
            \m{f}_m - \m{K}_{ms} \m{g}
        \end{bmatrix}
    \end{equation}

    where:
    \begin{eqarray}
        \m{\hat{K}}_{um} &= \m{K}_{um} + \m{K}_{us} \m{\closure{T}} \\
        \m{\hat{K}}_{mm} &= \m{K}_{mm} + \m{\closure{T}}^T \m{K}_{ms}^T
                            + \m{K}_{ms} \m{\closure{T}}
                            + \m{\closure{T}}^T \m{K}_{ss} \m{\closure{T}}
    \end{eqarray}

\end{bbox}


\subsubsection{MPC Application}

Afterwards the \textbf{MPCs} are applied:
\begin{equation}
    \m{u}_s = \m{\closure{T}} \m{u}_m
\end{equation}

This, when applied to the whole model:
\begin{equation}
    \begin{bmatrix}
        \m{u}_u  \\
        \m{u}_m \\
        \m{u}_s \\
    \end{bmatrix}
    = \begin{bmatrix}
        \m{I} & \m{0} \\
        \m{0} & \m{I} \\
        \m{0} & \m{\closure{T}} \\
    \end{bmatrix}
    \begin{bmatrix}
        \m{u}_u \\
        \m{u}_m \\
    \end{bmatrix}
\end{equation}

with \textbf{transformation matrix} $ \m{T} $ being:
\begin{equation}
    \m{T} = \begin{bmatrix}
        \m{I} & \m{0} \\
        \m{0} & \m{I} \\
        \m{0} & \m{\closure{T}} \\
    \end{bmatrix}
\end{equation}

Applying the transformation matrix to get \textbf{modified master equation system}:
\begin{eqarray}
    \m{T}^T \m{K}_g \m{T} \m{\hat{u}}_g &= \m{T}^T \m{f}_g \\
    \m{\hat{K}}_g \m{\hat{u}}_g &= \m{\hat{f}}_g
\end{eqarray}

Where $ \m{\hat{u}}_g = \begin{bmatrix} \m{u}_{g,u} & \m{u}_{g,m} \end{bmatrix}^T $.

From now on, the subscript $ _g $ (\textit{as global}) is implied.


\subsubsection{Rotated Nodal Basis}

When a node or an SPC is defined as a \textbf{rotated}/\textbf{skewed} or in
a different type of Coordinate system (e.g. cylindrical) a transformation matrix
needs to be applied to such nodes.

\textbf{Cartesian Coordinate system} definition by angle (2D):
\begin{equation}
    \m{u}_l = \m{T} \m{u}_g
\end{equation}


\subsubsection{SPC Application}

Then to solve the \textbf{master stiffness equations} for prescribed \textbf{BCs}
partitioned to \textbf{unconstrained} and \textbf{master} DOFs such as:
\begin{equation}
    \begin{bmatrix}
        \m{K}_{uu} & \m{\hat{K}}_{um} \\
        \m{\hat{K}}_{um}^T & \m{\hat{K}}_{mm} \\
    \end{bmatrix}
    \begin{bmatrix}
        \m{u}_u \\
        \m{u}_m \\
    \end{bmatrix}
    = \begin{bmatrix}
        \m{f}_{u} \\
        \m{f}_m \\
    \end{bmatrix}
\end{equation}

where:
\begin{eqarray}
    \m{\hat{K}}_{um} &= \m{K}_{um} + \m{K}_{us} \m{\closure{T}} \\
    \m{\hat{K}}_{mm} &= \m{K}_{mm} + \m{K}_{ms} \m{\closure{T}}
                      + (\m{K}_{ms} \m{\closure{T}})^T
                      + \m{\closure{T}}^T \m{K}_{ss} \m{\closure{T}}
\end{eqarray}

Any of these \textbf{DOFs} can be constrained, so this partitioning scheme looses
all purposes. Therefore the matrix is repartitioned again into
\textbf{independent} and \textbf{constrained} partitions:
\begin{equation}
    \begin{bmatrix}
        \m{K}_{ii} & \m{K}_{ic} \\
        \m{K}_{ic}^T & \m{K}_{cc} \\
    \end{bmatrix}
    \begin{bmatrix}
        \m{u}_i \\
        \m{u}_c \\
    \end{bmatrix}
    = \begin{bmatrix}
        \m{f}_{i} \\
        \m{f}_c \\
    \end{bmatrix}
\end{equation}


\subsubsection{Solving Linear Equations}

Where the \textbf{unknowns} are $ \m{u}_i $, $ \m{f}_c $:
\begin{equation}
    \begin{bmatrix}
        \m{K}_{ii} & \m{K}_{ic} \\
        \m{K}_{ic}^T & \m{K}_{cc} \\
    \end{bmatrix}
    \begin{bmatrix}
        \boxed{\m{u}_i} \\
        \m{u}_c \\
    \end{bmatrix}
    = \begin{bmatrix}
        \m{f}_i \\
        \boxed{\m{f}_c} \\
    \end{bmatrix}
\end{equation}

Unpacking the equations:
\begin{eqarray}
    \m{K}_{ii} \boxed{\m{u}_i} &+ \m{K}_{ic} \m{u}_c &= \m{f}_i \\
    \m{K}_{ic}^T \boxed{\m{u}_i} &+ \m{K}_{cc} \m{u}_c &= \boxed{\m{f}_c} \\
\end{eqarray}

Then first solve for the \textbf{unknown displacements}:
\begin{eqarray}
    \m{K}_{ii} \boxed{\m{u}_i} &= \m{f}_u - \m{K}_{ic} \m{u}_c \\
    \boxed{\m{u}_i} &= \m{K}_{ii}^{-1} \left( \m{f}_i - \m{K}_{ic} \m{u}_c \right) \\
\end{eqarray}

Lastly solve for the \textbf{unknown reaction forces}:
\begin{equation}
    \boxed{\m{f}_c} = \m{K}_{ii} \m{u}_i + \m{K}_{ic} \m{u}_c
\end{equation}


\subsubsection{Recovering Full Displacement Vector}

The full displacement vector is then recovered:
\begin{eqarray}
    \m{T} &= \begin{bmatrix}
        \m{I} & \m{0} \\
        \m{0} & \m{I} \\
        \m{0} & \m{\closure{T}} \\
    \end{bmatrix} \\
    \m{\hat{u}} &= \begin{bmatrix}
        \m{u}_i \\
        \m{u}_c \\
    \end{bmatrix}
    = \begin{bmatrix}
        \m{u}_u \\
        \m{u}_m \\
    \end{bmatrix} \\
    \m{u}_g &= \m{T} \m{\hat{u}} \\
    \begin{bmatrix}
        \m{u}_u \\
        \m{u}_m \\
        \m{u}_s \\
    \end{bmatrix}
    &= \begin{bmatrix}
        \m{I} & \m{0} \\
        \m{0} & \m{I} \\
        \m{0} & \m{\closure{T}} \\
    \end{bmatrix}
    \begin{bmatrix}
        \m{u}_u \\
        \m{u}_m \\
    \end{bmatrix}
\end{eqarray}




\newpage
\section{RBE3 -- Interpolation Coupling}

\subsection{Why RBE3 Is Different from RBE2}

An RBE2 answers the question: \textit{how must the surrounding nodes move if the
connection is perfectly rigid?} An RBE3 answers a different question:
\textit{what generalized motion of a reference point best represents the motion
of the surrounding nodes?}

This reversal is essential. The RBE3 does not require all attachment nodes to
follow one rigid plate. Their local deformation remains available. The reference
point is instead related to a weighted interpolation of those nodal motions.
Loads applied to the reference point are transferred back to the attachment
nodes by work equivalence.

Let the translational displacement of attachment node $i$ be $\m{u}_i$, its
position $\m{x}_i$ and its positive weight $w_i$. Define the weighted centroid

\begin{equation}\label{rbe3-weighted-centroid}
    \m{x}_c = \frac{\sum_i w_i\m{x}_i}{\sum_i w_i},
    \qquad
    \m{r}_i=\m{x}_i-\m{x}_c.
\end{equation}

For a rigid motion about that centroid one would have

\begin{equation}
    \m{u}_i \approx \m{u}_c + \m{\theta}\times\m{r}_i.
\end{equation}

The centroidal translation that minimizes the weighted displacement error is

\begin{equation}\label{rbe3-average-translation}
    \m{u}_c = \frac{\sum_i w_i\m{u}_i}{\sum_i w_i}.
\end{equation}

The corresponding least-squares rotation follows from

\begin{equation}\label{rbe3-rotation-normal-equations}
    \m{J}_w\m{\theta}
    = \sum_i w_i\,\m{r}_i\times(\m{u}_i-\m{u}_c),
\end{equation}

where

\begin{equation}
    \m{J}_w = \sum_i w_i
    \left[(\m{r}_i^T\m{r}_i)\m{I}-\m{r}_i\m{r}_i^T\right].
\end{equation}

The matrix $\m{J}_w$ has the same geometric form as an inertia tensor. Its
invertibility immediately exposes an important limitation: if the attachment
nodes do not span the directions required to identify a rotation, the
interpolation is singular. A collinear set of nodes, for example, cannot define
rotation about its own line from translational data alone.

If the reference point is at $\m{x}_r$, its interpolated translation is

\begin{equation}\label{rbe3-reference-displacement}
    \m{u}_r = \m{u}_c + \m{\theta}\times(\m{x}_r-\m{x}_c).
\end{equation}

Equations~\eqref{rbe3-average-translation}--\eqref{rbe3-reference-displacement}
may be collected into an interpolation operator

\begin{equation}\label{rbe3-interpolation-operator}
    \m{u}_r = \m{H}\m{u}_a,
\end{equation}

where $\m{u}_a$ collects the active attachment-node DOFs.

\subsection{Load Transfer from Virtual Work}

The cleanest derivation of RBE3 load distribution does not start from guessed
nodal forces. It starts from virtual work. For compatible virtual displacements,

\begin{equation}
    \delta\m{u}_r=\m{H}\delta\m{u}_a.
\end{equation}

Equating the virtual work of the reference resultant and the attachment-node
forces gives

\begin{equation}
    \m{f}_r^T\delta\m{u}_r
    = \m{f}_a^T\delta\m{u}_a.
\end{equation}

Therefore

\begin{bbox}
\begin{equation}\label{rbe3-force-transfer}
    \m{f}_a = \m{H}^T\m{f}_r.
\end{equation}
\end{bbox}

This transpose relation is the fundamental reason that an interpolation coupling
can both recover a reference motion and distribute a force or moment without
introducing an artificial material stiffness.

\subsection{A Simple Symmetric Example}

Consider four attachment nodes at the corners of a square of side $2a$, with the
reference point at its centre and equal weights. A pure reference force $F_z$
contains no moment about the centroid. By symmetry and
Eq.~\eqref{rbe3-force-transfer}, each corner receives

\begin{equation}
    F_{z,i}=\frac{F_z}{4}.
\end{equation}

A pure reference moment $M_z$ cannot be carried by a force at the centroid.
Instead it is represented by a self-equilibrated set of in-plane nodal forces.
Their resultant force is zero while their moment about the centroid is $M_z$.
This example illustrates the important distinction between \textit{resultant
equivalence} and a prescribed surface traction: many different nodal force
patterns can have the same total force and moment.

\subsection{RBE3 as a Deformable Distributed Support}

One of the most useful applications is the connection of a spring or bushing to
an extended deformable surface. Let an RBE3 reference point represent the average
motion of a flange or shell patch and let a spring connect that reference point
to ground. The spring stores the support energy,

\begin{equation}
    \Pi_s=\frac{1}{2}(\m{u}_r-\m{u}_g)^T\m{K}_s(\m{u}_r-\m{u}_g),
\end{equation}

while the RBE3 merely maps $\m{u}_r$ to the surface and distributes the reaction
by Eq.~\eqref{rbe3-force-transfer}. Local deformation modes of the attachment
surface that do not change the interpolated reference motion are not suppressed.
The support is therefore distributed without turning the attached patch into a
rigid plate.

A particularly useful construction employs coincident reference nodes. One node
belongs to the interpolation spider; another belongs to ground, another body or
a second spider. Translational and rotational springs, a bushing, or a set of
very stiff selected-direction springs between the coincident nodes can then
prescribe the desired support behaviour. Even when a selected generalized DOF is
made nearly rigid, the surface itself can still warp because the RBE3 does not
force every leg node to undergo one rigid-body motion.

\begin{bbox}
\textbf{Engineering note -- ``Deformable support'' does not mean a soft RBE3.}

The RBE3 has no support stiffness that is tuned to make it deformable. The
compliance comes from the structural surface and from the spring or bushing. The
RBE3 only defines how the generalized support motion and reaction are transferred.
\end{bbox}

\subsection{RBE3 Is Not Automatically a Pressure Load}

A common modelling shortcut is to connect an RBE3 to all nodes on a face and
apply a force at the reference point. The resulting nodal force field should not
be confused with a prescribed uniform pressure.

For a uniform pressure $p$ the consistent nodal resultant is governed by surface
integration. In a simple tributary-area interpretation,

\begin{equation}\label{rbe3-tributary-weight}
    F_i \approx p A_i,
\end{equation}

where $A_i$ is the surface area represented by node $i$. An interior node of a
regular mesh generally has a larger tributary area than an edge node, and a
corner node a smaller one still. A surface-based coupling that derives its
weights from area can therefore assign smaller nodal force to edge and corner
nodes than to interior nodes.

An explicitly defined RBE3 with equal nodal weights behaves differently: the
translational share associated with equal weights is equal, not area-consistent.
On a regular surface this over-represents the edge and corner nodes per unit
tributary area. Thus neither ``RBE3'' nor ``equal weights'' should be interpreted
as synonymous with ``uniform pressure''. If the traction distribution matters,
apply a true element surface load or construct weights from the appropriate
surface integration.

\begin{bbox}
\textbf{Common mistake -- Using an equal-weight RBE3 as a pressure replacement.}

The correct check is not only whether total force and moment are recovered. The
analyst must also decide whether the local traction distribution is physically
important. Saint-Venant-type arguments may make the difference irrelevant far
from the load introduction, but local stresses near the spider can be strongly
affected.
\end{bbox}

\subsection{Choosing the Spider Legs}

Connecting every available mesh node is rarely a goal in itself. In many
industrial models the spider legs are attached to representative corner,
perimeter or otherwise well-spaced nodes. There are several reasons.

\begin{enumerate}
    \item The selected nodes should span the translations and rotations that the
          coupling must transmit. Geometric coverage is more important than raw
          node count.
    \item With higher-order elements, corner nodes often already provide a
          stable low-order geometric skeleton. Adding every midside node may add
          many coefficients without adding an independent global motion.
    \item Large spiders create long constraint equations. They can increase
          assembly memory, matrix connectivity and solver work. The cost is
          particularly noticeable for interpolation constraints because many
          attachment DOFs contribute to each generalized reference DOF.
    \item Excessive or poorly chosen legs can interact with other constraints,
          contact definitions or neighbouring spiders and create redundant
          relations.
\end{enumerate}

There is no universal maximum number of legs. Some programs impose explicit
limits; others are limited only by memory or the internal representation of a
constraint equation. The practical rule is to use enough well-distributed nodes
to represent the intended load path and interpolation, not the largest set that
the preprocessor can select.

\begin{bbox}
\textbf{Engineering note -- More RBE3 legs do not monotonically improve a model.}

Convergence of the surrounding finite element mesh and convergence of the
coupling definition are separate questions. Once the spider adequately spans the
intended region, adding hundreds of additional legs may change little physically
while increasing assembly cost substantially.
\end{bbox}

\section{Fixed-Distance and Other Geometric Constraints}

A fixed-distance relation is fundamentally different from RBE1, RBE2 and RBE3.
For two points with current positions

\begin{equation}
    \m{x}_1=\m{X}_1+\m{u}_1,
    \qquad
    \m{x}_2=\m{X}_2+\m{u}_2,
\end{equation}

and prescribed distance $L$, a convenient exact constraint is

\begin{equation}\label{fixed-distance-constraint}
    g(\m{u})=\frac{1}{2}
    \left[(\m{x}_2-\m{x}_1)^T(\m{x}_2-\m{x}_1)-L^2\right]=0.
\end{equation}

Its linearization uses the current connecting vector
$\m{r}=\m{x}_2-\m{x}_1$,

\begin{equation}
    \delta g = \m{r}^T(\delta\m{u}_2-\delta\m{u}_1).
\end{equation}

Unlike the ordinary small-displacement RBE2 relation, the exact constraint is
nonlinear because its direction changes as the points rotate. A solver may
enforce Eq.~\eqref{fixed-distance-constraint} through Lagrange multipliers,
penalty or augmented Lagrangian techniques. A multiplier implementation is
especially transparent: the multiplier is the axial constraint force and the
constraint adds an extra unknown rather than eliminating a dependent DOF.

This is why a contact- or constraint-element implementation of a fixed distance
can behave numerically very differently from a classical RBE even if both are
used to make a connection appear ``rigid'' in a graphical model. The kinematic
relation and the enforcement technique must be identified separately.

Joints can be constructed in the same way by combining geometric constraints.
A revolute joint removes five relative rigid-body motions and leaves one rotation;
a spherical joint enforces point coincidence but leaves relative rotation; a
cylindrical joint leaves one translation and one rotation along a common axis.
Dedicated joint elements merely package these relations, often with nonlinear
updates, limits, friction or stiffness added.

\section{Remote Couplings and Reference Points}

Modern preprocessors commonly expose a \textit{remote point}, \textit{reference
point} or \textit{remote coupling} rather than asking the analyst to construct a
classical spider manually. Such an object normally consists of

\begin{enumerate}
    \item a geometric reference point with selected generalized DOFs,
    \item a region of mesh nodes to which it is coupled,
    \item a kinematic choice such as rigid or deformable/distributing, and
    \item an internal enforcement method chosen by the solver.
\end{enumerate}

A remote point is therefore a modelling abstraction, not a universal finite
element type. A rigid option commonly produces RBE2-like behaviour. A deformable
or distributing option commonly produces RBE3-like interpolation. A nonlinear
surface-based option may instead employ contact elements and an internal MPC
algorithm. The analyst should determine which behaviour is actually generated,
especially for finite rotations, thermal loading and modal analyses.

\subsection{Example Implementations in Commercial Solvers}

The following names are included only to help map the solver-independent concepts
onto familiar software. They are not a substitute for the respective solver
manuals.

\begin{center}
\begin{tabular}{|p{0.24\textwidth}|p{0.30\textwidth}|p{0.34\textwidth}|}
\hline
\textbf{Concept} & \textbf{ANSYS examples} & \textbf{PERMAS examples} \\
\hline
General linear constraint & \texttt{CE} & \texttt{MPC\_GENERAL} \\
\hline
Equal/coupled DOF & \texttt{CP}, \texttt{CE} & \texttt{MPC\_SAME}, \texttt{MPC\_ASSIGN} \\
\hline
Rigid coupling & \texttt{CERIG}, rigid remote coupling & \texttt{MPC\_RIGID} \\
\hline
Interpolation coupling & \texttt{RBE3}, deformable remote coupling & \texttt{MPC\_WLSCON} \\
\hline
Other specialized couplings & contact-based MPCs, joint/constraint elements & \texttt{MPC\_JOIN} and other specialized MPCs \\
\hline
\end{tabular}
\end{center}

The names illustrate a broader point: a mature solver usually contains several
constraint families because equal DOFs, rigid-body kinematics, weighted
least-squares interpolation and nonlinear geometric relations are not numerically
identical problems.

\section{Advanced Considerations}

\subsection{Small-Displacement and Finite-Rotation Formulations}

For an RBE2-type relation linearized in the initial configuration, the dependent
translation is commonly written

\begin{equation}\label{rbe-small-rotation-relation}
    \m{u}_s=\m{u}_m+\m{\theta}_m\times\m{r}_0,
\end{equation}

where $\m{r}_0$ is the initial arm. Equation~\eqref{rbe-small-rotation-relation}
is the tangent approximation to rigid-body motion. If a large rotation is imposed
while the equation remains unchanged, the slave point follows the initial tangent
rather than the exact circular trajectory.

A finite-rotation constraint must instead use the current rotation tensor,

\begin{equation}\label{rbe-finite-rotation-relation}
    \m{x}_s=\m{x}_m+\m{R}\m{r}_0,
\end{equation}

and be relinearized as the configuration changes. Small strain does not rescue a
small-rotation constraint: a structure may undergo tiny strains while rotating
through many degrees.

The distinction is formulation-specific. Turning on geometric nonlinearity in a
model does not guarantee that every legacy constraint object is upgraded to
Eq.~\eqref{rbe-finite-rotation-relation}. Some classical command-based RBE
implementations are explicitly restricted to small-deflection analysis, whereas
other contact-based or dedicated constraint elements update the relation in the
current configuration.

\subsection{Thermal Expansion of Rigid Couplings}

An ideal rigid spider defined only by mechanical kinematics preserves its initial
relative distances. During a uniform temperature change the surrounding material
may instead require

\begin{equation}
    L(T) \approx L_0\left(1+\alpha\Delta T\right).
\end{equation}

If the coupling itself does not admit this expansion, it can restrain an otherwise
free thermal strain and generate artificial reactions and stresses. Some rigid
or constraint-element implementations therefore permit a thermal expansion
coefficient or an equivalent temperature-dependent reference geometry.

The correct choice depends on what the rigid region represents. A truly rigid
insert with different thermal expansion may intentionally restrain the
surrounding structure. A spider used only as a numerical load-transfer device
usually should not create such an unintended thermal restraint.

\subsection{Modal and Dynamic Analyses}

An RBE2 can raise natural frequencies because it removes deformation modes from
the attached region. This is not an artificial effect if the physical connection
is genuinely rigid; it is artificial if the spider merely stands in for load
introduction to an otherwise flexible surface. RBE3 is often preferable in the
latter case because it preserves local structural deformation.

Constraint treatment also affects mass and inertia. An ideal MPC has no mass of
its own, but a reference node may carry lumped mass or rotary inertia, and
eliminated dependent DOFs transfer the structural mass matrix through the same
kinematic transformation used for stiffness. Dynamic models should therefore be
checked for both stiffness and mass consequences.

\subsection{Interaction with Contact and Other Constraints}

Overlapping rigid spiders, prescribed boundary conditions, bonded contact and
joint constraints can describe the same motion more than once. The result may be
redundant equations, singular multiplier blocks or unintended locking. A warning
about redundant constraints should therefore be treated as a modelling diagnostic,
not merely a solver nuisance.

RBE couplings should also not replace contact where separation or sliding is
physical. As emphasized in Chapter~\ref{chap:contact-formulations}, contact must
decide whether the constraint is active. An RBE remains a bilateral relation for
as long as it exists.

\section{Engineering Selection Guidelines}

A practical selection can be organized around a few questions.

\begin{enumerate}
    \item \textbf{Must the attached region itself remain rigid?}
          If yes, an RBE2-type rigid coupling is appropriate. If no, do not use a
          rigid spider merely to distribute a load.
    \item \textbf{Is the objective to transmit a resultant force, moment or
          generalized support motion while preserving local deformation?}
          An RBE3-type interpolation coupling is usually the natural starting
          point.
    \item \textbf{Does the connection possess finite compliance?}
          Model that compliance with a spring, bushing, beam, cohesive law or
          explicit geometry. An RBE3 may be used to distribute its endpoint
          reaction, but it does not replace the constitutive component.
    \item \textbf{Can the connection open, slide or change status?}
          Use contact or an appropriate nonlinear joint rather than a permanent
          RBE constraint.
    \item \textbf{Are large rotations expected?}
          Verify that the chosen implementation updates its constraint geometry.
          Do not infer finite-rotation capability from the RBE name alone.
    \item \textbf{Is thermal expansion important?}
          Check whether the coupling expands consistently with the represented
          physical part or introduces an unintended restraint.
    \item \textbf{Does the local traction distribution matter?}
          A resultant-equivalent RBE3 load is not automatically a uniform
          pressure. Use area-consistent weights or an actual surface load when
          necessary.
    \item \textbf{How many spider legs are really required?}
          Choose a well-spaced set that spans the intended motion. Avoid huge
          spiders unless their additional resolution has a demonstrated purpose.
\end{enumerate}

\begin{bbox}
\textbf{Final engineering rule.}

Choose the \textit{admissible relative motion} first, the \textit{constraint
formulation} second and the solver command last. A correct command applied to the
wrong physical constraint is still the wrong model.
\end{bbox}

\section{Summary}

Kinematic couplings are best understood as engineering definitions of admissible
relative motion. Chapter~\ref{chap:multifreedom-constraints} established the
numerical methods by which equality constraints can be enforced; this chapter
has focused on constructing the constraints themselves.

RBE1 represents a general rigid-coupling concept. RBE2 maps a reference-node
translation and rotation to dependent nodes and therefore creates a rigid region.
RBE3 reverses that logic: it interpolates generalized reference motion from a set
of attachment nodes and transfers forces back by virtual-work equivalence,
thereby preserving local deformation of the attached surface.

Fixed-distance and joint constraints show why the RBE family is only one part of
a broader subject. Their exact equations are geometric and generally nonlinear,
and they may be enforced by multipliers, penalty methods or other algorithms.
Contact, treated in Chapter~\ref{chap:contact-formulations}, extends the same
constraint viewpoint to inequalities and changing active sets.

Finally, the practical behaviour of a coupling depends on details that are easy
to overlook: small- versus finite-rotation kinematics, thermal expansion, local
load weighting, interaction with other constraints, and the number and placement
of spider legs. These details determine whether a convenient modelling shortcut
represents the intended engineering connection or silently changes the structure.
